\section*{第 \ref{ch:b1:ocean-circulation} 章——海洋环流}
% values: calc/ch29.py

\begin{solution}{exo:b1:ocean-circulation:1}
(a)~风速为 \qty{4}{m/s} 时（第一个公式）：$1000\,C_D = 0.29 + 0.775 + 0.481 = 1.546$，
$T = 1.3 \times \num{1.546e-3} \times 16 = \qty{0.0322}{Pa}$。风速为 \qty{7}{m/s}
时（第二个公式）：$1000\,C_D = 0.60 + 0.497 = 1.097$，$T = \qty{0.0699}{Pa}$。风速为
\qty{15}{m/s} 时：$1000\,C_D = 1.665$，$T = \qty{0.487}{Pa}$。(b)~$0.487/0.0699 =
6.97$，与 $(15/7)^2 = 4.59$ 相比：阻力系数本身随风速增大，
因为更强的风会掀起更粗糙的波浪，这些波浪更能阻滞气流。
\end{solution}

\begin{solution}{exo:b1:ocean-circulation:2}
(a)~北半球，向东吹的风右侧 \ang{90}：指向
南方。(b)~风向西南吹；其右侧为西北。
(c)~南半球，向东吹的风左侧 \ang{90}：指向
北方。(d)~风向北吹；其左侧为西，
在西海岸即为离岸方向。情形 (d) 中，输送指向离岸：表层
海水由下方补充，冷水沿海岸上升。
\end{solution}

\begin{solution}{exo:b1:ocean-circulation:3}
(a)~$31.7 + 3.5 - 17.9 = \qty{17.3}{\sverdrup}$，与已发表值
\qty{17.1}{\sverdrup} 相比：三部分平均值之和即平均\omterm{def:b1:ocean-circulation:moc}{经向翻转环流}
在舍入误差与不确定度（$\pm\qty{3.2}{\sverdrup}$）范围内相符。
(b)~在 \qty{1100}{m} 以上：$16.6 + 0.4 = \qty{17.0}{\sverdrup}$ 向北；以下：
$12.0 + 5.9 = \qty{17.9}{\sverdrup}$ 向南。两者几乎平衡，这是必然的，
在断面以北没有水体堆积的条件下；微小差异
位于不确定度范围内。(c)~$\num{17.1e6} \times \num{3.156e7} = \qty{5.40e14}{m^3}$，即
每年 \qty{5.40e5}{km^3}；$\num{1.3324e9}/(\num{5.40e5}) \approx 2470$ 年。
\end{solution}

\begin{solution}{exo:b1:ocean-circulation:4}
$D_E = 7.6 \times 10/\sqrt{\sin\varphi}$ 与 $V_0 = 0.0127 \times
10/\sqrt{\sin\varphi}$。在 \ang{15}：\qty{149}{m} 与 \qty{0.250}{m/s}
（为风速的 \qty{2.5}{\percent}）；在 \ang{45}：\qty{90}{m} 与 \qty{0.151}{m/s}
（\qty{1.5}{\percent}）；在 \ang{60}：\qty{82}{m} 与 \qty{0.136}{m/s}
（\qty{1.4}{\percent}）。在赤道附近 $\sin\varphi \to 0$，两个公式
都趋于发散；更根本地，$f \to 0$，使得用于平衡摩擦力所需的
科里奥利力趋于零，平衡建立所需的时间 $1/f$ 变得
长于风保持稳定的时间。埃克曼将其公式限制在 $|\varphi| \geq
\ang{10}$ 的范围内。
\end{solution}

\begin{solution}{exo:b1:ocean-circulation:5}
(a)~$f = 2\Omega\sin\ang{11} = \qty{2.783e-5}{s^{-1}}$；$M = 0.07/(\num{2.783e-5}) =
\qty{2515}{kg.m^{-1}.s^{-1}}$，向西吹的风右侧 \ang{90}：
向北。(b)~$Q = (2515/1027) \times \num{5e6} = \qty{1.22e7}{m^3/s}$，
\qty{12.2}{\sverdrup}。(c)~该估计位于直接
测量（$12.0 \pm 5.5$）的不确定度范围内，且接近由长期风场得到的
值（$13.5 \pm 0.3$）：该定理解释了观测到的输送。
\end{solution}

\begin{solution}{exo:b1:ocean-circulation:6}
(a)~$f(\ang{2}) = \qty{5.09e-6}{s^{-1}}$；$M = 0.05/(\num{5.09e-6}) =
\qty{9824}{kg.m^{-1}.s^{-1}}$，\qty{9.57}{m^2/s}：在北纬 \ang{2} 处指向北方
（向西吹的风的右侧）；在南纬 \ang{2} 处指向南方（左侧）。(b)~沿\omterm{def:b1:spherical-astronomy:geographic}{经度}方向长
\qty{1}{m} 的条带内，$2 \times 9.57 = \qty{19.1}{m^3/s}$ 穿越
相距 \qty{444.8}{km} 的两条纬线流出：$w = 19.1/(\num{4.448e5}) = \qty{4.30e-5}{m/s}$，
每天约 \qty{3.7}{m}。(c)~在赤道处 $f = 0$，$|\vec M| = T/|f|$ 为
无穷大：此处不存在稳态埃克曼平衡，上述数值对所选\omterm{def:b1:spherical-astronomy:geographic}{纬度}非常敏感，
只是数量级上的估计。可保留的只有
符号：赤道两侧的输送都背离赤道，因此表层
海水辐散，冷水沿之上升。
\end{solution}

\begin{solution}{exo:b1:ocean-circulation:7}
(a)~$f(\ang{15}) = \qty{3.775e-5}{s^{-1}}$：$M = 0.06/(\num{3.775e-5}) =
\qty{1590}{kg.m^{-1}.s^{-1}}$，指向北方；$f(\ang{40}) =
\qty{9.375e-5}{s^{-1}}$：$M = \qty{1280}{kg.m^{-1}.s^{-1}}$，指向南方。
(b)~两条纬线之间（相距 \qty{2780}{km}），$\partial M_y/\partial y =
-(1590 + 1280)/(\num{2.78e6}) = \qty{-1.03e-3}{kg.m^{-2}.s^{-1}}$：由
\cref{prop:b1:ocean-circulation:pumping}，$w_E = \qty{-1.005e-6}{m/s}$，向下的
速度为每年 \qty{31.7}{m}。(c)~$\num{1.005e-6} \times \num{2.78e6} \times
\num{5e6} = \qty{1.40e7}{m^3/s}$，\qty{14.0}{\sverdrup} 向下输送。
\end{solution}

\begin{solution}{exo:b1:ocean-circulation:8}
(a)~$H = 1027 \times 3992 \times \num{17e6} \times 12 = \qty{8.36e14}{W}$，
\qty{0.84}{PW}。(b)~它是实测 \qty{1.2}{PW} 的 \qty{70}{\percent}：真实的
上层分支在表层附近比 \qty{15}{\celsius} 的平均值更暖，且
环流也输送热量。(c)~$1027 \times 3992 \times \num{17e6} \times 3 =
\qty{2.1e14}{W}$，\qty{0.21}{PW}，为 (a) 的四分之一。只有在返回前变冷的水
才携带大量热量：大西洋的热输送依赖于
深层水的形成。
\end{solution}

\begin{solution}{exo:b1:ocean-circulation:9}
设 $x$、$y$ 与 $1 - x - y$ 分别为 AAIW、NADW 和 AABW 的占比。盐与热量
平衡，以 AABW 为基准：
\[
  -0.52\,x + 0.33\,y = 34.731 - 34.67 = 0.061,\qquad
  2.9\,x + 3.6\,y = 2.78 - 0.4 = 2.38 .
\]
行列式为 $-0.52 \times 3.6 - 0.33 \times 2.9 = -2.829$，由克莱姆法则
给出 $x = (0.061 \times 3.6 - 0.33 \times 2.38)/(-2.829) = 0.20$，$y = (-0.52
\times 2.38 - 2.9 \times 0.061)/(-2.829) = 0.50$。该样本为 \qty{20}{\percent}
AAIW、\qty{50}{\percent} NADW 与 \qty{30}{\percent} AABW，
这是三种\omterm{def:b1:fluid-envelopes-structure:ts-diagram}{水团}所构成三角形内的一点。
\end{solution}

\begin{solution}{exo:b1:ocean-circulation:10}
(a)~令 $W = u + iv = V_0e^{i\pi/4}e^{(1+i)az}$，其实部与虚部即为
命题中的公式，$W'' = (1+i)^2a^2W = 2ia^2W$。两个方程
合并为 $A_z(u'' + iv'') = -fv + ifu = if(u + iv)$，即 $A_zW'' = ifW$，
这成立，因为 $2ia^2A_z = if$ 对 $a^2 = f/(2A_z)$ 成立。(b)~$W'(0) =
(1+i)aV_0e^{i\pi/4} = \sqrt2\,aV_0\,e^{i\pi/2} = i\sqrt2\,aV_0$：表层应力
$\rho A_zW'(0)$ 指向北方，大小为 $\rho A_z\sqrt2\,aV_0 = \rho A_z\sqrt2
\sqrt{f/(2A_z)}\;T/(\rho\sqrt{fA_z}) = T$。(c)~$\int_{-\infty}^0 W\,\mathrm dz =
V_0e^{i\pi/4}/((1+i)a) = V_0/(\sqrt2\,a)$，为实数：输送没有
向北分量，其向东分量为 $\rho V_0/(\sqrt2\,a) =
T/\sqrt{fA_z}\cdot\sqrt{A_z/f} = T/f$。这就是该定理：大小为 $T/f$，
方向在向北吹的风右侧 \ang{90}。
\end{solution}

\begin{solution}{exo:b1:ocean-circulation:11}
(a)~沿北向取一小步 $\mathrm dy = a_\oplus\,\mathrm d\varphi$，则 $f = 2\Omega
\sin\varphi$ 改变 $2\Omega\cos\varphi\,\mathrm d\varphi$，于是 $\beta = 2\Omega
\cos\varphi/a_\oplus = 2 \times \num{7.292e-5} \times \cos\ang{30}/(\num{6.371e6}) =
\qty{1.98e-11}{m^{-1}.s^{-1}}$。(b)~$R/\beta = \num{e-6}/(\num{1.98e-11}) =
\qty{5.04e4}{m}$，约 \qty{50}{km}。(c)~内部：$\num{30e6}/(\num{5e6} \times
1000) = \qty{6.0e-3}{m/s}$，即 \qty{6}{mm/s}；边界流：$\num{30e6}/(\num{5.04e4}
\times 1000) = \qty{0.60}{m/s}$，快了一百倍，即宽度之比。墨西哥湾流
的速度为 \qty{1.8}{m/s}，更快：实际洋流比
这一估计更窄，其流动集中在表层附近，且其再循环
又增加了输送量。
\end{solution}

\begin{solution}{exo:b1:ocean-circulation:12}
(a)~凡海深超过 \qty{1000}{m} 之处，\qty{1000}{m} 以上的水柱
至多含 \qty{1}{km} 的水，其他处则更少：
\qty{1000}{m} 以上的体积至多为 $\num{361.84e6} \times 1 = \qty{3.618e8}{km^3}$，于是
以下的体积至少为 $\num{1.3324e9} - \num{3.618e8} = \qty{9.71e8}{km^3}$。
(b)~$\num{26e6} \times \num{3.156e7} = \qty{8.21e14}{m^3}$，每年
即 \qty{8.21e5}{km^3}；$\num{9.71e8}/(\num{8.21e5}) \approx 1200$ 年，与
传送带的千年量级相同。(c)~$(1/2)^{25/12.3} = 0.244$，约四分之一；
$\num{e-3} \times 25 \times \num{3.156e7} = \qty{7.9e5}{m}$，约 \qty{790}{km}。
25 年后，氚扩散的距离不超过缓慢漂移所能携带的范围，
而其约四分之一仍可探测到。
\end{solution}

\begin{solution}{pb:b1:ocean-circulation:1}
\begin{enumerate}
  \item 北风沿海岸向南吹。
  \item $1000\,C_D = 0.60 + 0.071 \times 8 = 1.168$：$C_D = \num{1.168e-3}$。
  \item $T = 1.3 \times \num{1.168e-3} \times 64 = \qty{0.0972}{Pa}$。
  \item $C_D = \num{1.452e-3}$，$T = \qty{0.272}{Pa}$：应力增大为 $2.80$ 倍，
        超过 $(12/8)^2 = 2.25$，因为 $C_D$ 本身随风速增大。
  \item $\Omega = \qty{7.292e-5}{rad.s^{-1}}$；$f = 2\Omega\sin\ang{38} =
        \qty{8.979e-5}{s^{-1}}$。
  \item $2\pi/f = \qty{7.00e4}{s}$，\qty{19.4}{h}。科里奥利力使一股洋流
        在此时间内转过一整圈；在风吹了
        约这么长时间之前，水体做振荡，而不是
        趋于定理所述的稳态平衡。
  \item 在向南吹的风右侧 \ang{90}：
        指向西方，即外海。
  \item $M = T/f = 0.0972/(\num{8.979e-5}) = \qty{1082}{kg.m^{-1}.s^{-1}}$。
  \item $1082/1027 = \qty{1.054}{m^2/s}$，对应每米海岸。
  \item $D_E = 7.6 \times 8/\sqrt{\sin\ang{38}} = \qty{77.5}{m}$。
  \item $1.054/77.5 = \qty{0.0136}{m/s}$，即 \qty{1.36}{cm/s}。
  \item $V_0 = 0.0127 \times 8/\sqrt{\sin\ang{38}} = \qty{0.129}{m/s}$，在风右侧 \ang{45}，
        指向西南。漂流浮标测得的值约为
        埃克曼值的一半，约 \qty{0.07}{m/s}。
  \item \omterm{def:b1:ocean-circulation:ekman}{埃克曼层}中离开海岸的水，无法由海岸另一侧
        补充，也无法由上方补充：近岸处由该层以下
        上升的水补充，如
        \cref{sec:b1:ocean-circulation:upwelling}中剖面图所示；等温线向岸边
        上翘。
  \item $1.054 \times 1000 \times 86400 = \qty{9.11e7}{m^3}$，每天每千米
        海岸。
  \item $w = 1.054/(\num{2e4}) = \qty{5.27e-5}{m/s}$，即每天 \qty{4.55}{m}。
  \item $77.5/4.55 = 17$ 天。
  \item 冷水自下方冷却空气：近地表的
        气层变冷且稳定，水汽可能凝结成雾和低层云，
        因而无对流，海岸降雨很少。
  \item 上升的水富含营养物质，它们滋养浮游植物，
        浮游植物又滋养浮游动物和鱼类。
  \item 向北吹的风驱动输送，方向在其右侧 \ang{90}，
        指向东：表层水堆积在海岸边并下沉（\omterm{def:b1:ocean-circulation:upwelling}{下降流}）；
        \omterm{def:b1:ocean-circulation:upwelling}{上升流}停止，较暖的表层水重新回到海岸。
  \item $f(\ang{20}) = \qty{4.988e-5}{s^{-1}}$：$M = \qty{1948}{kg.m^{-1}.s^{-1}}$，
        $\sin\ang{38}/\sin\ang{20} = 1.80$ 倍。
  \item 在南半球，输送方向位于
        风左侧 \ang{90}。要使其向西离岸，风必须向北吹，即
        从南方吹来：两种情形下，风都沿海岸向赤道方向吹。
  \item $1000\,C_D = 0.29 + 0.62 + 0.308 = 1.218$；$T = \qty{0.0396}{Pa}$；$M =
        \qty{441}{kg.m^{-1}.s^{-1}}$；两者之比均为 $0.407$。
  \item 沿海岸每米，$M c_p\,\Delta T = 1082 \times 3992 \times 6 =
        \qty{2.59e7}{W}$，即 \qty{25.9}{MW}；沿 \qty{1000}{km} 海岸则为 \qty{2.59e13}{W}，
        即 \qty{25.9}{TW}。
  \item $\num{31.7e6}/1.054 = \qty{3.01e7}{m}$，约 \qty{30000}{km}：\omterm{def:b1:ocean-circulation:upwelling}{沿岸
        上升流}输送的水远少于\omterm{def:b1:ocean-circulation:wbc}{西边界流}。
  \item $Q = 1.054 \times \num{e6} = \qty{1.05e6}{m^3/s}$：\omterm{def:b1:ocean-circulation:upwelling}{沿岸上升流}的输送量
        沿 \qty{1000}{km} 海岸为 \qty{1.05}{\sverdrup}。
\end{enumerate}
\end{solution}
